\subsection{The complete ordinary five-state budget at order thirteen}

\begin{lemma}[Proved: the three-by-four module dimensions]
\label{CP2:A-five-dimensional-three-four-target}
Let $A$ be a five-dimensional unital rational algebra, $V$ a
three-dimensional left module and $W$ a four-dimensional right
module. They cannot satisfy
\[
 \dim\operatorname{Tor}^A_j(W,V)=
 \begin{cases}1,&j=0,2,\\0,&j=1\text{ or }j\geq3.\end{cases}
\]
Consequently a monoid with minimal band $4\times5$ and five
outside elements cannot have type $\CP^2$.
\end{lemma}
\begin{proof}
If $A$ is local, Lemma~\ref{A:local-target-facts} makes its
residue field rational and both modules cyclic. The cyclic
cover $A\to W$ has one-dimensional kernel, necessarily the
right simple module $\mathbb Q$. Since $V$ is not projective,
dimension shifting gives
\[
 \operatorname{Tor}_3^A(W,V)
   \cong\operatorname{Tor}_2^A(\mathbb Q,V)\ne0.
\]
This contradicts the proposed pattern.

For nonlocal $A$, Lemma~\ref{A:nonlocal-five-target} leaves
the product of a four-dimensional local algebra with
$\mathbb Q$, the zero-composite three-vertex line algebra,
and the five-dimensional two-vertex Nakayama algebra.
For $A=B\times\mathbb Q$, positive Tor comes entirely from
the $B$ summands. They have nonzero tensor product by the
local top quotient, so that tensor product must have dimension
one. They are both cyclic over $B$. A cyclic module of
dimension four is free; one of dimension three has simple
first syzygy; and one of dimension one is simple. In those
cases either positive Tor vanishes or the required vanishing
in degrees one and three is impossible by the local minimal
resolution argument. The only remaining dimensions are two
on both sides, which are excluded by
Lemma~\ref{A:local-four-small-target}.

We use the complete indecomposable lists and displayed
projective resolutions proved in
Proposition~\ref{A:three-by-six-CP2}. For the line algebra
write $L_a=Ae_2$, $L_b=Ae_3$, $R_a=e_1A$, $R_b=e_2A$.
The indecomposables on the left are
$S_1,S_2,S_3,L_a,L_b$, of dimensions $1,1,1,2,2$;
the right list is $T_1,T_2,T_3,R_a,R_b$ with the same
dimensions. Nonzero degree-two Tor of dimension one forces
exactly one $S_3$ and one $T_1$. Degree-one vanishing then
forbids $S_2$ and $T_2$. The complete remaining possibilities
and their tensor dimensions are
\[
\begin{array}{c|ccc}
 &T_1+3T_3&T_1+T_3+R_a&T_1+T_3+R_b\\\hline
 2S_1+S_3&5&5&3\\
 S_3+L_a&3&2&2\\
 S_3+L_b&6&2&3.
\end{array}
\]
No entry is one. Here plus signs in module expressions
denote direct sums.

For the Nakayama algebra the left list is
$S_1,S_2,L,Ae_1,Ae_2$, of dimensions $1,1,2,3,2$;
the right list is $T_1,T_2,R,e_1A,e_2A$ with the same
dimensions. Let their multiplicities be
$(a,b,c,d,e)$ and $(A_1,B_1,C_1,D_1,E_1)$, respectively.
The exact resolutions give
\[
 \dim\operatorname{Tor}_2=(B_1+C_1)(b+c),\qquad
 \dim\operatorname{Tor}_1
   =A_1(b+c)+B_1a+C_1(a+c).
\]
The required dimensions force $b+c=B_1+C_1=1$.
If $c=1$, the left dimension forces $a=1$, which makes
the displayed first Tor positive. If $b=1$, the remaining
left dimension is either $a=2$ or $e=1$; the first again
makes first Tor positive. Thus $V=S_2\oplus Ae_2$ and
$A_1=0$. The four-dimensional right module must be either
$T_2\oplus e_1A$ or $R\oplus e_2A$. Direct evaluation of
the tensor quotient gives dimension three in both cases.
This excludes the last alternative.

The exact successor
\path{monoid/code/five_dimensional_three_four_pairs.py}
reconstructs both literal algebra multiplications, every
indecomposable action and every displayed resolution, and
enumerates all $480$ and $144$ module pairs of the required
dimensions. It finds no Tor triple $(1,0,1)$. Its four-rule
derivation preserves the earlier checker, changing only the
requested dimensions, target assertion and output. The full
Tor tables and resolutions also agree entry for entry with
the previously checked two-by-five control. The source,
derivation and completed result are bound under
\path{two_object19/ordinary13_five_algebra/}.

For the monoid assertion, the band--Tor modules have exactly
these dimensions, and the contracted algebra has dimension
five. A projective plane would require the excluded pattern.
\end{proof}

\begin{lemma}[Proved: full core-orbit quotients retain a finite model]
\label{CP2:A-full-column-orbit-wedge}
Suppose $M$ has rectangular minimal band $I\times J$ and
every core orbit of a finite nonempty right act $S$ has
exactly $|J|$ states. If there are $o$ such orbits and the
original act cone $X$ is simply connected, the actual
constant-action quotient map extends to a homotopy equivalence
\[
 X\vee\bigvee^{o-1}S^2\ \simeq\
 Y=B\mathcal C(M,\bar S).
\]
In particular $Y$ has finite CW homotopy type when $X$ does.
\end{lemma}
\begin{proof}
Lemma~\ref{CP2:A-full-column-act-kernel} gives the exact
integral right-module sequence with kernel
$(e\mathbb ZM)^{\oplus(o-1)}$. In addition to its homology
isomorphisms in degrees at least three, the long exact Tor
sequence gives
\[
 0\longrightarrow H_2(X;\mathbb Z)
 \longrightarrow H_2(Y;\mathbb Z)
 \longrightarrow\mathbb Z^{o-1}
 \longrightarrow H_1(X;\mathbb Z)
 \longrightarrow H_1(Y;\mathbb Z)\longrightarrow0,
\]
because $e\mathbb ZM\otimes_{\mathbb ZM}\mathbb Z
\cong\mathbb Z$. All maps are those of the actual quotient.

Constant-action compression identifies $Y$ with the
classifying space of a monoid whose minimal ideal is
$I\times\bar S$. Such a monoid has trivial universal group:
every band element is idempotent and hence maps to the
group identity, while $me$ is a band element for every
monoid element $m$ and any band idempotent $e$. It follows
that $Y$ is simply connected.

Since $X$ is simply connected, the displayed sequence has
free cokernel $\mathbb Z^{o-1}$ in degree two. Choose a
splitting. The Hurewicz isomorphism for $Y$ represents a
basis of its complementary summand by based two-spheres.
Wedge these maps with the actual quotient map from $X$.
The resulting map is an integral homology isomorphism
between simply connected CW complexes, so homological
Whitehead makes it a homotopy equivalence. This proves
both the claimed wedge and its compatibility with the
original map; no finiteness or boundedness is assumed
merely from a second-cohomology comparison.
\end{proof}

\begin{proposition}[Proved: four rows, five outside elements, five states]
\label{CP2:A-four-row-five-outside-five-state}
An ordinary right-act cone with source minimal band
$4\times2$, exactly five outside elements, and at most
five states cannot have type $\CP^2$.
\end{proposition}
\begin{proof}
The source has order thirteen. At most four states are
covered by Theorem~\ref{CP2:A-four-state-fourteen}, so suppose
$|S|=5$ and the cone has projective-plane type. A singleton
core orbit gives a global fixed state and zero second
products. Thus every core orbit is a full two-point orbit.
There are either one or two such orbits.

With one core orbit, the band acts constantly on $S$.
Actual compression gives a monoid of the same homotopy
type with minimal band $4\times5$ and five outside
elements, contradicting the first lemma.

With two core orbits, the constant-action quotient
identifies corresponding columns and leaves the one
transient state distinct. Hence $|\bar S|=3$. The preceding
lemma and compression give
\[
 BN\simeq\CP^2\vee S^2,
 \qquad K_N\cong4\times3,
 \qquad\dim_{\mathbb Q}H^2(BN;\mathbb Q)=2.
\]
This is an actual finite CW homotopy type. The fully proved
rank-two theorem
Theorem~\ref{CP2:A-three-by-four-rank-two}, including its
opposite orientation, says all rational second products
on $BN$ vanish. This contradicts the nonzero square on
the projective-plane summand. Its finite-CW hypothesis
has just been established, not transferred through an
uncontrolled actor quotient.
\end{proof}

\begin{theorem}[Proved: the nineteen-arrow ordinary-act budget]
\label{CP2:A-oriented-nineteen-thirteen-one-five}
No ordinary right-act cone over a monoid of order at most
thirteen with at most five states has type $\CP^2$.
In particular the complete nineteen-arrow budget
$(13,1,5)$ is excluded. The left-act version holds by
opposition.
\end{theorem}
\begin{proof}
Proposition~\ref{CP2:A-thirteen-five-state-single-shape}
reduces a hypothetical example to the precise shape
excluded by the preceding proposition. This covers all
minimal groups and all nonfaithful source actions without
enumerating original thirteen-element monoids.
\end{proof}

\begin{remark*}[Separate opposite-corner question]
The proof establishes the actual ordinary-act assertion.
The budget $(13,2,4)$ remains separate: a rational orbit
comparison by itself does not transfer the actual monoid
minimum or the finite-CW hypotheses used here.
\end{remark*}
